% =============================================================================
% Lecture 5: Physical Data
% Physical AI: Machine Learning Systems That Sense and Act (Volume IV), Chapter 5
% =============================================================================
% Thesis: a robot's own actions choose the states it will meet, so a
% demonstration archive is engineered for recovery coverage and for the
% meaning, timing, and origin of every label.
% Running case: the spring-latched cage-door latch campaign (Napkin Math 5.1
% to 5.3; sections 5.7, 5.8, 5.11). Hook and payoff: DROID's recalibration of
% about 36,000 episodes (section 5.10.1).
%
% Build (from slides/): make deck V=4 D=05_data; figures: images/make_figs.py;
% every derived number: images/calcs.py. Plan: lecture_map.yml.
% =============================================================================

\documentclass[aspectratio=169, 12pt]{beamer}
\usepackage{../../assets/beamerthememlsys}
\mlsyssetup{
  volume={Volume IV},
  chapter={Chapter 5},
  chaptertitle={Physical Data},
  layer=world,
  logo={../../assets/img/logo-mlsysbook.png},
  instlogo={../../assets/img/logo-harvard.png},
}
\usepackage[T1]{fontenc}
\usepackage[scaled=0.9]{helvet}
\usepackage{courier}
\renewcommand{\familydefault}{\sfdefault}
\usepackage{booktabs,amsmath,amssymb}
\usetikzlibrary{arrows.meta,positioning,calc}
\graphicspath{{images/}}
\setcounter{mlsystotalsections}{7}

\title{Physical Data}
\author{Vijay Janapa Reddi}
\institute{Harvard University}
\date{}

% A large answer choice for predict and vote slides
\newcommand{\choice}[2]{\textbf{\textcolor{mlsysaccent}{#1}}\enspace #2}
% A goals or takeaways line with a text-free thumbnail
\newcommand{\thumbline}[2]{%
  \begin{minipage}[c]{2.6cm}\centering\includegraphics[width=2.5cm,height=0.95cm,keepaspectratio]{#1}\end{minipage}%
  \hspace{0.3cm}\begin{minipage}[c]{0.74\textwidth}\raggedright\small #2\end{minipage}}
% Verbatim-like artifact text, small
\newcommand{\art}[1]{{\fontsize{7.6}{9}\selectfont\ttfamily #1}}

\begin{document}

% =============================================================================
% TITLE
% =============================================================================
\mlsysnote{[0.5 min] [CORE] Read the thesis aloud once. This lecture opens Part II, Teaching the Machine. The notes total about 71 minutes; 12.5 are tagged [OPTIONAL], and cutting them leaves a complete lecture of about 58.5 minutes. The running case is the latch campaign: the warehouse mobile manipulator opening a spring-latched cage door. The chapter is marked for a second lecture later; the natural split falls after the Labels section and the yield slides (Physical Data, \S5.6).}
\mlsystitle{Physical Data}{A robot's own actions choose the states it will meet, so a demonstration archive is engineered for recovery coverage and for the meaning, timing, and origin of every label.}{cover_data_blueprint.png}

% =============================================================================
% SECTION: THE ARCHIVE
% =============================================================================
\mlsysnote{[0.5 min] [CORE] Open on a real dataset that passed every file check and was still wrong.}
\mlsysdivider{Why can a large archive of successful demonstrations leave a robot unable to recover from its first mistake?}{Physical Data, \S5.1 and \S5.10.1}
\section{The Archive}

\begin{frame}{DROID recalibrated 36,000 episodes that parsed without error}
  \centering
  \safeimg[width=0.86\textwidth,height=5.0cm,keepaspectratio]{droid-calibration-grid.jpg}\\[0.1cm]
  {\small 76,000 trajectories, 350 hours, 564 scenes, one robot model.}
  \mlsyscredit{Khazatsky et al., DROID, RSS 2024, Figure 13 (CC BY 4.0), via Physical Data, figure 9. Scale: Physical Data, table 5. Recalibration: \S5.10.1 (Lu et al. 2025; Khazatsky et al. 2024).}
  \note{[2.5 min] [CORE] NARRATE: DROID is one of the largest open robot-manipulation datasets: 76,000 trajectories, 350 hours, 564 scenes, collected by many labs on one robot model. Each tile is a real frame from a DROID camera with a purple silhouette of the arm drawn on top, computed from the recorded joint angles and the camera's stored calibration. Where the silhouette sits on the arm, the camera pose is right. An independent re-analysis found that rendered silhouettes barely overlapped the visible arm across many episodes. The DROID team measured the overlap, discarded failed sessions, re-estimated camera poses, and published new calibrations for about 36,000 episodes. Every one of those files had loaded without error. ENGAGE: hold the question for the next slide. WARN: the grid shows the corrected calibrations; the failure was in the original ones. FLEX: keep this slide; the lecture returns to it in the Provenance section.}
\end{frame}

\begin{frame}{What was wrong with those episodes?}
  \mlsysactivity{Predict}{1 min}
  \centering
  {\normalsize The arrays had the right shapes, types, and ranges. What was wrong?}\\[0.5cm]
  \begin{minipage}{0.8\textwidth}\large
  \choice{A}{Corrupted video frames}\\[0.2cm]
  \choice{B}{Joint angles out of range}\\[0.2cm]
  \choice{C}{Stored camera poses that did not match the cameras}\\[0.2cm]
  \choice{D}{Missing episodes}
  \end{minipage}
  \mlsyscredit{Physical Data, \S5.10.1.}
  \note{[1 min] [CORE] ENGAGE: students vote with fingers. Do not reveal. The answer is C, and the lecture returns to it at the end. A and D are what a parser would catch; B is what a range check would catch. C passes both, because a wrong camera pose is a well-formed matrix. WARN: students who know DROID may say ``bad labels''; ask what kind of label, and keep the question open. FLEX: if short on time, ask the question aloud from the previous slide.}
\end{frame}

\begin{frame}{A demonstration costs an operator, a workcell, and time}
  \centering
  \safeimg[width=0.88\textwidth,height=4.6cm,keepaspectratio]{aloha2-fleet.jpg}\\[0.1cm]
  {\small One person per workcell, one demonstration at a time.}
  \mlsyscredit{ALOHA 2 Team (Google DeepMind), ALOHA 2, arXiv:2405.02292, Figure 1 (top), cropped.}
  \note{[1 min] [OPTIONAL] NARRATE: a real data-collection fleet, Google DeepMind's ALOHA 2 workcells. Each station has a pair of leader arms that an operator moves by hand and a pair of follower arms that copy them, with cameras overhead and on the wrists. The paper calls the fleet capable of thousands of demonstrations per day. Compare this with how a language model's corpus is gathered: no one sits in a chair for each sentence. LINK: in Lecture 1 the causal boundary was where a command becomes physical work; every demonstration here crosses it, so every demonstration costs operator time and component life. ENGAGE: what does one more demonstration cost here that one more web page does not? (Expected: an operator's time, wear on the arms, and a reset of the scene.) FLEX: skip if short; the yield slides price this later.}
\end{frame}

\begin{frame}{Recall: what may the permission path change?}
  \mlsysactivity{Retrieval}{1 min}
  \centering
  {\Large\bfseries Close your notes.}\\[0.25cm]
  {\normalsize Who may write to the drives? What can happen to a proposal\\before it reaches them, and what does the record of that hold?}\\[0.4cm]
  \uncover<2->{%
    \thumbline{thumb_kit_stack_nervous.pdf}{Only the permission path writes to the drives. It may admit, clamp, or refuse a proposal, and the handoff record dates each step.}}
  \mlsyscredit{Lecture 4 (sole-writer actuation authority; handoff record). Lecture 1 (proposal boundary).}
  \note{[1 min] [CORE] Retrieval from Lecture 4. Give 40 seconds. Step 2: the permission path is the sole writer to the drives; a proposal from the learned side may be admitted unchanged, clamped, or refused, and the handoff record dates each of those events on one clock. LINK: Lecture 1's proposal boundary separates the side that may only propose from the side that may command. Today this matters because every demonstration we record passed through that path. A logged command may be the operator's request or the clamped command that actually moved the arm, and those are different labels.}
\end{frame}

\begin{frame}{Why can perfect demonstrations leave a robot stuck?}
  \normalsize
  By the end of today, you will be able to:\\[0.25cm]
  \thumbline{thumb_corridor_step2.pdf}{Explain why a policy's first error leads to states its demonstrations never covered.}\\[0.15cm]
  \thumbline{thumb_taps_step4.pdf}{Choose which of four logged signals a dataset imitates, and date each pair on one clock.}\\[0.15cm]
  \thumbline{thumb_scenario_step2.pdf}{Price a collection campaign in cell time and wear, and audit coverage against a declared table.}\\[0.15cm]
  \thumbline{thumb_takeover_step3.pdf}{Curate a takeover and write the provenance record that travels with an archive.}\\[0.2cm]
  {\footnotesize\textcolor{midgray}{Thread: the second law, endogenous data. Record: the provenance record.}}
  \note{[1 min] [CORE] The title is the driving question. One line per segment pair: recovery, labels, cost and coverage, interventions and provenance. LINK: Lecture 1 named the four laws. This lecture opens Part II and states the second, endogenous data, the thread named on the bottom line: a machine's own actions choose the states, and therefore the data, it collects. Unlike the first law, irreversibility, it is a design principle the book elevates; only the first law is physical. The record today adds to the record chain is the provenance record, the third record after Lecture 2's limit record and Lecture 4's handoff record. Running case: the warehouse mobile manipulator opening a spring-latched cage door, with a 15 N contact tripwire and two approach speeds, a guarded 0.03 m/s and a faster 0.10 m/s.}
\end{frame}

% =============================================================================
% SECTION: RECOVERY
% =============================================================================
\mlsysnote{[0.5 min] [CORE] Start with why a robot's data differs from an image dataset at all.}
\mlsysdivider{Why do more clean demonstrations not teach recovery?}{Physical Data, \S5.1 to \S5.3}
\section{Recovery}

\begin{frame}{A robot's actions choose the states it sees next}
  \centering
  \mlsyssteps[width=0.95\textwidth,height=4.9cm,keepaspectratio]{loops_step}{2}\\[0.1cm]
  {\small\uncover<2->{\textbf{Endogenous covariate shift:} the policy's own actions lead it to states its data never held.}}
  \mlsyscredit{Physical Data, \S5.1, \S5.2, and definition 5.1.}
  \note{[2 min] [CORE] Step 1: an image classifier. A wrong label on image 1 does not change image 2; the test set is fixed before the model runs. Covariate shift in ordinary machine learning means the inputs change for outside reasons, such as daylight to night. Step 2: a robot policy. Its action moves the body, which moves the camera and changes contact, so the next observation depends on the last action. A small error therefore changes what the policy sees next. The chapter calls this endogenous covariate shift, endogenous because the robot causes it. ENGAGE: name one exogenous shift and one endogenous shift for the latch task. (Expected: afternoon sunlight on the door; a gripper that missed the centerline by a millimeter and now sees the latch from a new angle.) WARN: students often treat closed-loop shift as bad luck in the test set; it is caused by the policy itself, so collecting more of the same data cannot remove it.}
\end{frame}

\begin{frame}{Clean demonstrations fill the corridor and nothing beside it}
  \centering
  \mlsyssteps[width=0.95\textwidth,height=5.5cm,keepaspectratio]{corridor_step}{2}
  \mlsyscredit{Physical Data, figure 1 and \S5.3, redrawn at slide size.}
  \note{[2 min] [CORE] Top view of the latch. Step 1: every retained demonstration drives the gripper forward along a narrow corridor centered on the strike plate. Skilled operators rarely drift, and runs that drifted were aborted or thrown away during curation, so the archive holds thousands of forward actions. Step 2: the gripper sits 3 mm toward the fixture. The archive holds no example that starts here and returns to the corridor, so nothing constrains the policy's output at this state. Whether the drift dies out depends on the plant. Compliant joints may pull the gripper back; a stiff arm keeps drifting until the permission path or a person stops it. ENGAGE: why would curation make this worse? (Expected: discarding runs that drifted removes the only off-corridor states the archive had.) WARN: students blame noisy labels; the problem is missing actions at missing states.}
\end{frame}

\begin{frame}{ALVINN met the same gap in 1989 and synthesized recoveries}
  \begin{columns}[c]
    \begin{column}{0.56\textwidth}
      \centering\safeimg[width=\textwidth,height=4.6cm,keepaspectratio]{navlab-1-5.jpg}
    \end{column}
    \begin{column}{0.40\textwidth}
      \small
      Drivers kept the van centered, so the images held almost no drift.\\[0.3cm]
      Pomerleau shifted each image sideways and corrected the steering label.\\[0.3cm]
      A gripper off-axis also feels new contact forces; no image warp creates those.
    \end{column}
  \end{columns}
  \mlsyscredit{Photo: Carnegie Mellon Navlab 1 (farthest) to 5 (front), Firefly4342, CC BY-SA 4.0, Wikimedia Commons. Physical Data, \S5.3 footnote 1 (Pomerleau 1989, 1991).}
  \note{[1 min] [OPTIONAL] NARRATE: Carnegie Mellon's Navlab vehicles, 1 at the back through 5 at the front. ALVINN, a neural network that steered from camera images, drove the early Navlabs. It learned from a human driver who kept the van centered, so its training images held almost no off-center views, the corridor problem of the last slide. Pomerleau's fix shifted each camera image sideways in software and adjusted the steering label to match, inventing recovery examples about 20 years before DAgger. ENGAGE: why does that trick not transfer to the latch? (Expected: an off-center road is just a different viewpoint of a flat surface; an off-axis gripper also experiences different contact forces, which no image transform can produce.) FLEX: cut if short; the point is that the recovery gap is decades old.}
\end{frame}

\begin{frame}{Cloning's worst-case cost grows as the horizon squared}
  \centering
  \mlsyssteps[width=0.88\textwidth,height=4.5cm,keepaspectratio]{bound_step}{2}\\[0.05cm]
  {\small \textbf{Behavioral cloning}: supervised learning of the demonstrator's action at each recorded state.}
  \mlsyscredit{Physical Data, \S5.3 and equation 1 (Ross and Bagnell 2010); comparison at $S = 500$ steps (5 s at 100 Hz), $\varepsilon = 0.001$, \S5.8. Curves derived.}
  \note{[2.5 min] [CORE] Behavioral cloning fits the demonstrator's action at each recorded state, so its loss is measured only where the demonstrator went. Let $\varepsilon$ bound the chance of a wrong action per step and $S$ be the episode length in control steps. Step 1: if errors did not feed back, the expected cost would grow as $S\varepsilon$: 0.5 at 500 steps. Step 2: in closed loop, a first error at step $t$ can leave the policy in unsupported states for the remaining $S - t$ steps. Summing over $t$ gives about $S^2\varepsilon/2$; the published bound (Ross and Bagnell 2010) is $S^2\varepsilon$, 250 at 500 steps, half the 500-step maximum. Above $\varepsilon = 1/S = 0.002$ the bound exceeds the episode maximum and says nothing. ENGAGE: at $S = 1{,}000$ steps, above what $\varepsilon$ is the bound vacuous? (Expected: 0.001.) WARN: this bounds task cost, position error needs a bound of its own, and this one is a worst case; with log loss or chunked, multimodal policies the scaling differs (Foster et al. 2024; Simchowitz et al. 2025).}
\end{frame}

\begin{frame}{Will 1,000 more hours of clean openings fix it?}
  \mlsysactivity{Vote}{2 min}
  \begin{minipage}{0.9\textwidth}\normalsize
  \choice{A}{Yes. More data always lowers the error.}\\[0.15cm]
  \choice{B}{No. Support beside the corridor stays at zero.}\\[0.15cm]
  \choice{C}{Only with a larger model.}\\[0.15cm]
  \choice{D}{Only if the logger runs faster than 1 kHz.}
  \end{minipage}\\[0.4cm]
  \uncover<2->{\centering\normalsize B.\enspace An hour at 1 kHz: \textbf{3,600,000} records, about \textbf{1,800} independent samples.}
  \mlsyscredit{Physical Data, \S5.3 (nominal demonstration trap) and equation 2: $N(1-\rho)/(1+\rho)$ with $\rho = 0.999$ (derived).}
  \note{[2 min] [CORE] ENGAGE: vote, then discuss with a neighbor for 30 seconds. Step 2: B. A thousand more hours of centered openings add density along the corridor; the number of demonstrations beside it stays zero. C does not help, because a larger model only fits the same states better. D is the volume illusion. A logger at 1 kHz writes 3.6 million state-action records an hour, but adjacent records 1 ms apart differ by sensor noise. For a slowly varying statistic with lag-one correlation 0.999, the effective sample size $N(1-\rho)/(1+\rho)$ is about 1,800, three orders of magnitude fewer. A faster clock oversamples the same trajectory and never reaches an unvisited state. WARN: the effective-sample-size formula is for estimating one mean from an autocorrelated series; the coverage argument does not depend on it. FLEX: drop the correlation numbers if short.}
\end{frame}

\begin{frame}{Most public LeRobot datasets hold fewer than 16 episodes}
  \centering
  \safeimg[width=0.92\textwidth,height=4.9cm,keepaspectratio]{lerobot_hub.pdf}
  \mlsyscredit{Hugging Face Hub API, datasets tagged LeRobot, retrieved 2026-10-12 (MLSysIM Literature.LeRobotHubMonthly, LeRobotEpisodeHistogram); Physical Data, figure 5. Earlier counts are lower bounds.}
  \note{[1.5 min] [CORE] Real data from the Hugging Face Hub, where the open-source LeRobot toolchain publishes recorded sessions. Left: the count of public LeRobot-format datasets, 1,085 at the end of 2024, 26,548 a year later, and 80,132 in October 2026, about 74 times growth. Right: episodes per dataset for the 59,031 datasets whose card records a count. 58 percent hold fewer than 16 episodes; only 2 percent hold 512 or more; the median falls in the 8-to-15 bin. ENGAGE: is community data growing wide or deep? (Expected: wide; many robots, scenes, and operators, few repeats of any one machine's states.) Land the segment: the archive misleads because the robot's own errors lead it where the demonstrations never went; more of the same data, from one rig or from a hub, does not fill that gap. WARN: earlier counts are lower bounds, since deleted datasets are not counted.}
\end{frame}

% =============================================================================
% SECTION: LABELS
% =============================================================================
\mlsysnote{[0.5 min] [CORE] The data comes from a rig. What, exactly, does each recorded action mean?}
\mlsysdivider{What does a recorded action mean, and when was it true?}{Physical Data, \S5.4 to \S5.6.1}
\section{Labels}

\begin{frame}{Leader arms copy joint angles to the follower directly}
  \centering
  \safeimg[width=0.80\textwidth,height=4.6cm,keepaspectratio]{gello-teleop.jpg}\\[0.1cm]
  {\small Seven collection paths, each with its own artifact in the recording.}
  \mlsyscredit{Wu et al., GELLO, IROS 2024, arXiv:2309.13037, Figure 1 (two UR5s, an xArm, a Franka). Physical Data, \S5.4 and table 1.}
  \note{[1.5 min] [CORE] Real rigs: GELLO leader devices built for two UR5 arms (A), an xArm (B), and a Franka (C). The operator moves a small leader arm whose joints mirror the follower's, and the follower copies the joint angles, with no inverse kinematics in between. The chapter names seven collection paths, each with its own artifact in the recording: kinesthetic teaching hides the cameras behind the demonstrator's body; teleoperation adds display delay and hesitation; scripts repeat one rigid path; exploration risks collisions; deployment logs are censored by the safety layer; hand-held grippers record no robot joints; simulation records the solver's contact model. ENGAGE: which path gives the archive recovery data most cheaply? (Expected: none for free; exploration finds boundaries but risks hardware; teleoperation with deliberate perturbation is the usual answer.) WARN: a leader arm feels the leader's own friction; the follower's contact force never reaches the operator's hand.}
\end{frame}

\begin{frame}{One command exists as four logged signals}
  \centering
  \mlsyssteps[width=0.95\textwidth,height=5.0cm,keepaspectratio]{taps_step}{4}\\[0.05cm]
  {\small\uncover<4->{\textbf{Action-label semantics}: the tap a dataset keeps decides what the policy learns.}}
  \mlsyscredit{Physical Data, \S5.5.1 and figure 2, redrawn. PTP: IEEE 1588 Precision Time Protocol, a hardware clock shared across processors.}
  \note{[3 min] [CORE] One operator motion becomes four numbers. Step 1, Asked: the operator's or policy's raw request on the application processor, before any check. Step 2, Mapped: the request turned into a joint setpoint by kinematic mapping, still on the proposal side. Step 3, Allowed: the command after the permission path at 1 kHz has checked and possibly clamped it; this is what actually reached the drives. Step 4, Done: what the encoders and current sensors measured. All four go into one log stamped by one hardware clock, here IEEE 1588 PTP. Which tap is the label decides what the policy learns. For the latch, the label is the mapped command when the permission path passed it unchanged; a modified command stays an attributed intervention. LINK: Lecture 4's handoff record dates these same events. ENGAGE: when are taps 2 and 3 equal? (Expected: whenever the permission path admitted the command unchanged.) WARN: tap 4 is the next state, an input to the policy, never the action label.}
\end{frame}

\begin{frame}{Which tap did each policy imitate?}
  \mlsysactivity{Diagnose}{2 min}
  \small
  \renewcommand{\arraystretch}{1.15}
  \begin{tabular}{@{}>{\raggedright\arraybackslash}p{0.52\textwidth}>{\raggedright\arraybackslash}p{0.40\textwidth}@{}}
    \toprule
    \textbf{Symptom at deployment} & \leavevmode\uncover<2->{\textbf{Tap it imitated}} \\
    \midrule
    Repeats fast pushes the machine always clamped. & \leavevmode\uncover<2->{\textcolor{RedLine}{\textbf{1 Asked}}: moves the plant never ran.} \\
    Stops against the latch as if stopping were the goal. & \leavevmode\uncover<2->{\textcolor{RedLine}{\textbf{4 Done}}: zero velocity while holding torque.} \\
    Chatters at 8 to 12 Hz near the plate. & \leavevmode\uncover<2->{Raw operator velocity: hand tremor copied.} \\
    \bottomrule
  \end{tabular}
  \mlsyscredit{Physical Data, \S5.4.1 and \S5.5.1; Checkpoint 5.1.}
  \note{[2 min] [CORE] ENGAGE: pairs name the tap behind each symptom, two minutes. Step 2: Row 1 trained on the request, tap 1: the operator asked for more than the permission path allowed, and the transition followed the clamped command, so the policy learned a move the machine refuses. Row 2 trained on the measured response, tap 4: against the latch the measured velocity is zero while the motors hold torque, so ``stop'' looks like the command. Row 3 trained on raw human velocities: physiological tremor at 8 to 12 Hz rides on the motion and the policy copies it as torque chatter, which wears the harmonic drives. WARN: none of these is a network architecture problem; each is a labeling decision made before training. FLEX: if short, do rows 1 and 2 only.}
\end{frame}

\begin{frame}{DROID and LeRobot log one command and the measured state}
  \begin{columns}[T]
    \begin{column}{0.56\textwidth}
      {\footnotesize\textbf{\textcolor{mlsysaccent}{DROID RLDS schema, per step (excerpt)}}}\\[0.1cm]
      \art{"observation": \{\\
      \ \ "joint\_position": tf.Tensor(7, ...),\\
      \ \ "cartesian\_position": tf.Tensor(6, ...),\\
      "action\_dict": \{\\
      \ \ "joint\_position": ...,\ \# commanded\\
      \ \ "cartesian\_velocity": ..., \# commanded\\
      "reward": ..., \# 1 on final step for demos}
    \end{column}
    \begin{column}{0.38\textwidth}
      {\footnotesize\textbf{\textcolor{mlsysaccent}{LeRobot ALOHA info.json}}}\\[0.1cm]
      \art{"robot\_type": "aloha",\\
      "total\_episodes": 50,\\
      "fps": 50,\\
      "observation.state": [14],\\
      "observation.effort": [14],\\
      "action": [14]}
    \end{column}
  \end{columns}
  \vspace{0.35cm}
  \centering{\small\uncover<2->{Neither records what a safety layer changed between command and motion.}}
  \mlsyscredit{DROID Docs, Dataset Schema, droid-dataset.github.io; Hugging Face, lerobot/aloha\_static\_coffee meta/info.json. Retrieved 2026-10-19; verbatim, ``...'' marks cuts.}
  \note{[1.5 min] [OPTIONAL] Two real schemas, copied from the source files. Left: DROID's per-step fields. The observation holds the measured joint and Cartesian state, tap 4; the action dictionary holds commanded positions and velocities in several forms, which corresponds to taps 1 or 2. Right: a real LeRobot dataset recorded on an ALOHA rig: 50 episodes at 50 frames per second, 14 measured joint states, 14 efforts, and one 14-element action. Step 2: neither schema has a field for an enforced command, tap 3, or for a clamp or override flag. For tabletop research rigs without an independent safety layer that is reasonable; for a machine with a permission path it loses the evidence of every intervention. ENGAGE: in which field would a clamp show up? (Expected: nowhere; it would appear only as a gap between action and next state.) WARN: do not read this as a criticism of these datasets; they record what their rigs had.}
\end{frame}

\begin{frame}{A 35 ms pairing lag fits the guarded approach only}
  \centering
  \mlsyssteps[width=0.88\textwidth,height=4.6cm,keepaspectratio]{pairing_step}{2}\\[0.05cm]
  {\small Pairing error $\approx$ speed $\times$ lag. Date each frame at its \textbf{exposure midpoint}.}
  \mlsyscredit{Physical Data, Napkin Math 5.1, equation 3, and figure 3 (illustrative speeds and clearance; ceilings derived). Rolling-shutter skew 15 to 33 ms: \S5.5.2.}
  \note{[2 min] [CORE] A logger that stamps a camera frame when it arrives in memory pairs each action with an image taken earlier; here 35 ms earlier. The recorded state is off by speed times lag. Step 1: at a 0.50 m/s transit, 17.5 mm, 8.75 times the 2 mm latch clearance; at the guarded 0.03 m/s approach, 1.05 mm, inside it; at the faster 0.10 m/s, 3.5 mm, outside it. Step 2: the lag ceilings are clearance over speed: 66.7 ms guarded, 20 ms fast. The gray band is a second error. A rolling-shutter camera exposes rows one after another, 15 to 33 ms top to bottom, so a correctly stamped frame can still break the 20 ms ceiling. The fix is a hardware trigger and a timestamp at the exposure midpoint. LINK: Lecture 1's stopping budget turned delay into distance the same way; Lecture 8 owns the capture timestamp. WARN: this is a logging error, removable; the operator's own delay on the next slides is not.}
\end{frame}

\begin{frame}{How much lag fits 1 mm at 0.10 m/s?}
  \mlsysactivity{Estimate}{1 min}
  \centering
  {\normalsize A tighter strike plate leaves 1 mm of clearance.\\The fast approach is 0.10 m/s. The logger lags 35 ms.}\\[0.5cm]
  \uncover<2->{{\Large $1\ \text{mm} \div 0.10\ \text{m/s} = \mathbf{10\ ms}$}\\[0.3cm]
  {\normalsize The 35 ms logger fails by 3.5$\times$; only the guarded approach could use it.}}
  \mlsyscredit{Physical Data, Checkpoint 5.1 (derived).}
  \note{[1 min] [CORE] ENGAGE: 45 seconds, individually. Step 2: the lag ceiling is clearance divided by speed: 1 mm divided by 0.10 m/s, which is 0.1 mm per ms, gives 10 ms. A 35 ms logger produces 3.5 mm of pairing error, 3.5 times the clearance. At the guarded 0.03 m/s it produces 1.05 mm, just over the tighter 1 mm, so even the guarded approach now fails by a hair. WARN: students divide speed by clearance and get a frequency; check units aloud (mm divided by mm per ms). FLEX: skip the guarded follow-up if time is short.}
\end{frame}

\begin{frame}{The operator's reaction takes 200 of the loop's 340 ms}
  \begin{columns}[c]
    \begin{column}{0.70\textwidth}
      \centering\safeimg[width=\textwidth,height=2.6cm,keepaspectratio]{teleop_strip.pdf}\\[0.3cm]
      {\small Keep both times in every pair: the frame's exposure midpoint and the command's issue time.}
    \end{column}
    \begin{column}{0.27\textwidth}
      \centering\safeimg[width=\textwidth,height=3.4cm,keepaspectratio]{argonne-master-slave.jpg}\\[0.05cm]
      {\footnotesize Goertz, 1949: force came back through the linkage.}
    \end{column}
  \end{columns}
  \mlsyscredit{Physical Data, \S5.6.1 and equation 5; illustrative setup (human share 58.8\%, derived). Photo: Raymond Goertz and his master-slave manipulator, 1949, Argonne National Laboratory, CC BY-SA 4.0, Wikimedia Commons.}
  \note{[1.5 min] [OPTIONAL] The teleoperation loop in the chapter's illustrative setup runs as follows. The display shows a contact 50 ms after it happens; the operator perceives, decides, and moves in 200 ms; the command reaches the follower in 10 ms; the arm takes 80 ms to respond. 340 ms in all, 58.8 percent of it the human. So a command issued at $t + 250$ ms answers a state the operator saw at $t$. Unlike the logger's lag, this delay is a real property of the demonstration, and no timestamp correction removes it. The record keeps both times and the session's measured loop latency, and training either conditions on enough history or pairs each command with the frame the operator was answering. UMI matches the same latency at deployment. Right: Raymond Goertz in 1949 with the master-slave manipulator he built at Argonne; its linkage returned contact force to the hand, which a video-only rig drops. ENGAGE: why do operators slow down before contact? (Expected: a fast correction loop around 340 ms of delay would oscillate.) WARN: the operator's delay varies per command, so no single fixed shift recovers it.}
\end{frame}

\begin{frame}{57 percent of pooled episodes were logged at 10 Hz}
  \centering
  \safeimg[width=0.88\textwidth,height=4.8cm,keepaspectratio]{oxe_rate.pdf}
  \mlsyscredit{Open X-Embodiment dataset overview sheet, 72 datasets, 2,419,193 episodes (MLSysIM Literature.OpenXEmbodimentComposition, verified 2026-10-12); Physical Data, figure 8(b). Chunk rate: Lecture 3.}
  \note{[1 min] [CORE] Real data: the Open X-Embodiment collection, the largest pooled robot corpus, weighted by episodes and grouped by the command rate each dataset declares. 57 percent of its 2.4 million episodes were logged at 10 Hz; 28 percent declare no rate at all, often because their actions are primitive skills; only 2 percent were recorded at or above 20 Hz, the running machine's chunk renewal rate from Lecture 3. A label's timing is part of its meaning, and most of this corpus cannot be retimed to our machine without a declared rate. Land the segment: a label means one of four signals, dated at one instant on one clock, and a record without the tap, the clock, or the rate cannot be interpreted later. FLEX: if short, state the 57 and 28 percent and move on.}
\end{frame}

% =============================================================================
% SECTION: COVERAGE
% =============================================================================
\mlsysnote{[0.5 min] [CORE] Natural split for a second lecture falls here (end of \S5.6). Now price the data and ask what it covers.}
\mlsysdivider{What does an hour of collection buy, and which states does it cover?}{Physical Data, \S5.6.2 to \S5.7}
\section{Coverage}

\begin{frame}{One scheduled hour keeps 14.7 minutes of demonstrations}
  \centering
  \mlsyssteps[width=0.90\textwidth,height=3.4cm,keepaspectratio]{yield_step}{2}\\[0.2cm]
  {\normalsize\uncover<2->{52 attempts $\times$ 0.68 kept $=$ \textbf{35.36} demonstrations, \textbf{24.6\%} of the hour.}}\\[0.1cm]
  {\small\uncover<2->{5,000 kept demonstrations: 7,353 attempts, \textbf{141.4 h} of cell time.}}
  \mlsyscredit{Physical Data, Napkin Math 5.3, equation 6, and table 3 (illustrative session; derived).}
  \note{[2.5 min] [CORE] The latch campaign's illustrative hour. Step 1: 8 minutes of startup and bus synchronization, 7 of sensor calibration, 10 of operator rest; the remaining 35 minutes hold 52 attempt-and-reset cycles of 25 s plus 15 s, so 13 minutes of resets and about 22 of attempts. Step 2: 20 percent of attempts fail in the trial (latch bind, the 15 N tripwire) and 15 percent of the rest fail quality filtering, so 0.68 survive: 35.36 demonstrations, 14.73 minutes, 24.6 percent of the scheduled hour. The campaign target of 5,000 kept demonstrations needs 7,353 attempts and 141.4 hours of cell time, 4.07 times the 34.7 hours of motion it keeps. ENGAGE: a faster door reset halves the 15 s. Which bill shrinks? (Expected: cell time only; wear is per attempt.) WARN: these session values are illustrative; teams must measure their own.}
\end{frame}

\begin{frame}{A 5,000-demonstration campaign uses 29.4\% of the cable's life}
  \begin{columns}[c]
    \begin{column}{0.46\textwidth}
      \centering\safeimg[width=\textwidth,height=4.3cm,keepaspectratio]{strain-wave-gear.jpg}\\[0.05cm]
      {\footnotesize Strain-wave gear: 14 reversals per attempt.}
    \end{column}
    \begin{column}{0.50\textwidth}
      \centering\safeimg[width=\textwidth,height=4.4cm,keepaspectratio]{wear.pdf}
    \end{column}
  \end{columns}
  \mlsyscredit{Photo: strain-wave gear set (Harmonic Drive AG), Pieceofmetalwork, CC BY-SA 4.0, Wikimedia Commons. Physical Data, Napkin Math 5.3 (illustrative ratings; derived).}
  \note{[1.5 min] [CORE] Left: a real strain-wave gear set, the wrist reducer type the chapter assumes: a flexible toothed cup inside a rigid ring, driven by an elliptical wave generator. Right: wear for the 5,000-demonstration campaign. Each attempt, kept or not, reverses torque on the wrist gear 14 times and bends the internal cable bundle 8 times. 7,353 attempts give 102,942 gear cycles, 2.1 percent of a 5-million-cycle rating, and 58,824 cable bends, 29.4 percent of a 200,000-bend rating. Three campaigns would use 88 percent of the cable's rated life, and a cracked conductor corrupts the sensor streams being recorded. LINK: Lecture 2's thermal limits and limit record bound the same body during collection. ENGAGE: why does the retention rate divide both bills? (Expected: failed attempts wear parts and cost time but are not kept.) WARN: a rated life is a population statistic (the life 90 percent of parts reach); one cable may fail sooner.}
\end{frame}

\begin{frame}{Raw logging writes 3.12 GB/s; the taps are 0.007\% of it}
  \centering
  \safeimg[width=0.90\textwidth,height=4.4cm,keepaspectratio]{bytes.pdf}
  \mlsyscredit{Physical Data, Napkin Math 5.2 (camera geometry datasheet; logging formats chosen; derived). SATA 550 and NVMe Gen4 7,000 MB/s: chapter values.}
  \note{[1 min] [OPTIONAL] One session logging every camera raw. The 4056 by 3040 navigation camera at 60 Hz alone writes 2,219 MB/s, 4.0 times what a SATA SSD sustains. All streams together, 3,122 MB/s, fit an NVMe Gen4 drive but fill 11.24 TB per hour. Eightfold compression brings the session to 390 MB/s at the cost of artifacts in the training frames. The kinematic stream that carries all four action taps for nine axes at 1 kHz is 0.22 MB/s, 0.007 percent of the bytes. ENGAGE: under buffer pressure, which stream must a logger drop last? (Expected: the taps; they are the cheapest and the only record of what the permission path did.)}
\end{frame}

\begin{frame}{Two archives of a million records cover 8 and 510 cells}
  \centering
  \safeimg[width=0.90\textwidth,height=4.5cm,keepaspectratio]{grids.pdf}\\[0.05cm]
  {\small The narrower archive has the lower held-out error.}
  \mlsyscredit{Physical Data, \S5.7 and table 4 (illustrative contact task: $\pm$15$^\circ$ at 1$^\circ$ bins, 0 to 20 N at 1 N bins, 600 cells).}
  \note{[2 min] [CORE] An illustrative contact task: contact angle from minus to plus 15 degrees and normal force up to 20 N, binned at one degree and one newton, 600 cells. Two collectors each record one million transitions at 100 Hz. The cautious operator stays within 2 degrees of perpendicular at 4 to 6 N: 8 cells, 1.3 percent. An exploratory policy covers the full angle range from 1 to 18 N: 510 cells, 85 percent. Yet the cautious archive has the lower held-out action error, 0.001 against 0.012 rad, because a narrow, repetitive archive is easy to fit. ENGAGE: which number would a dashboard show by default? (Expected: record count and validation loss, the two that cannot see the difference.) WARN: wide coverage does not prove the exploratory labels are good actions; only closed-loop evaluation (Lecture 7) tests competence. FLEX: one line for the curse of dimensionality: a full 28-dimensional state grid at 10 bins per axis has $10^{28}$ cells, beyond any fleet.}
\end{frame}


\begin{frame}{The latch archive fills 1 of the 6 cells the door requires}
  \centering
  \mlsyssteps[width=0.92\textwidth,height=4.75cm,keepaspectratio]{scenario_step}{3}\\[0.05cm]
  {\small\textbf{Operational design domain (ODD)}: required conditions, declared before collection.}
  \mlsyscredit{Physical Data, \S5.7, definition 5.3, and figure 6; latch-recovery-02 is a stipulated teaching branch (\S5.8).}
  \note{[2.5 min] [CORE] The operational design domain, ODD, is declared from what deployment requires, before any data is collected. A scenario table cuts it into a few task parameters; here plate offset and approach speed. Coverage is a join of where the archive went against that table. Step 1: against an uncritical table holding only the nominal plate at the guarded speed, the archive looks complete. Step 2: against the table the door actually needs, three offsets by two speeds, it fills one cell; five cells are known absences. The telemetry did not change; the standard of audit did. Step 3: the stipulated recovery collection latch-recovery-02, at +3 mm and 0.03 m/s, fills one more; four absences remain. ENGAGE: why must the ODD come before the data? (Expected: an ODD written from where the robot already succeeded certifies itself.) WARN: an occupied cell supplies training targets; it does not show the policy succeeds there. LINK: Lecture 16's release manifest carries these absences.}
\end{frame}

\begin{frame}{Single arms supply 91\% of Open X-Embodiment's episodes}
  \centering
  \safeimg[width=0.88\textwidth,height=4.6cm,keepaspectratio]{oxe_morph.pdf}
  \mlsyscredit{Open X-Embodiment dataset overview sheet (MLSysIM Literature.OpenXEmbodimentComposition, verified 2026-10-12); Physical Data, figure 8(a).}
  \note{[1 min] [OPTIONAL] Real data: the same 2.4 million pooled episodes by robot morphology, on a log axis. 91 percent come from single arms and 89 percent from tabletop scenes. Mobile manipulators, the running machine's class, contribute 119,386 episodes, 5 percent, and two datasets supply 95 percent of those. ENGAGE: our latch table has six cells; how many of these 2.4 million episodes are in our cells? (Expected: none are known to be; coverage is defined against our task, and the pool cannot report it.) The pool can help pretraining; it cannot fill the latch table.}
\end{frame}

\begin{frame}{Add a third speed. How many absences now?}
  \mlsysactivity{Estimate}{1 min}
  \centering
  {\normalsize The table adds an approach speed between 0.03 and 0.10 m/s.\\How many cells, and how many are known absences for the original archive?}\\[0.5cm]
  \uncover<2->{{\Large $3 \times 3 = \mathbf{9}$ cells;\enspace $\mathbf{8}$ known absences}\\[0.3cm]
  {\normalsize The archive did not change. The audit did.}}
  \mlsyscredit{Physical Data, Checkpoint 5.2 (derived).}
  \note{[1 min] [CORE] ENGAGE: 30 seconds. Step 2: three offsets times three speeds is 9 cells; the original archive still fills only the nominal, guarded cell, so 8 are known absences. Land the segment: physical data is bought in cell time and component life, and what it covers is a join against a table declared from the ODD, never a property the archive can report about itself. WARN: students add one cell instead of a column of three; draw the third column. FLEX: if behind, state the answer and move to the divider.}
\end{frame}

% =============================================================================
% SECTION: INTERVENTIONS
% =============================================================================
\mlsysnote{[0.5 min] [CORE] Absences get filled from the robot's own failures, but only after curation.}
\mlsysdivider{What should an archive keep from a takeover?}{Physical Data, \S5.8}
\section{Interventions}

\begin{frame}{Four different causes leave the same zero count}
  \small
  \renewcommand{\arraystretch}{1.45}
  \begin{tabular}{@{}>{\raggedright\arraybackslash}p{0.44\textwidth}>{\raggedright\arraybackslash}p{0.50\textwidth}@{}}
    \toprule
    \textbf{Why the +3 mm cell holds zero} & \textbf{What to do} \\
    \midrule
    Never visited & Collect there on purpose. \\
    Permission path cut the motion short (15 N tripwire) & Leave it to the permission path, which owns the stop. \\
    A person took over first & Curate: the drift states are the missing recovery data. \\
    Script discarded the failed run & Pipeline defect; recover the raw log if it exists. \\
    \bottomrule
  \end{tabular}
  \mlsyscredit{Physical Data, \S5.8 (four mechanisms of absence).}
  \note{[1.5 min] [CORE] A zero in the scenario table does not say why it is zero. Row 1: the collector never went there; targeted collection fixes it. Row 2: the permission path intercepted the motion, as the 15 N tripwire does on a misaligned latch; that state belongs to the safety layer, so do not teach the policy to go there. Row 3: a supervisor disengaged the policy before contact; the states the learner visited before takeover are exactly the off-corridor states the archive lacks. Row 4: the robot entered the state and failed, and an ingestion script discarded the episode as corrupt; the archive cannot show this at all. ENGAGE: which row can the dataset itself never reveal? (Expected: row 4; the evidence was deleted before serialization.) WARN: treating rows 2 and 3 alike either trains the policy to imitate an abort or throws away the boundary states it needs.}
\end{frame}

\begin{frame}{Curation splits a takeover into four intervals}
  \centering
  \mlsyssteps[width=0.88\textwidth,height=4.9cm,keepaspectratio]{takeover_step}{3}
  \mlsyscredit{Physical Data, \S5.8, equation 7, and figure 7(b), redrawn; 5 s lookback from the 4.47 s takeover budget, rounded up. Tracking error is schematic.}
  \note{[3 min] [CORE] One takeover on the latch. Step 1: the learner drifts from $t_{\mathrm{div}}$ until a supervisor takes over at $t_{\mathrm{take}}$. The divergence comes first, by at least the supervisor's response time; the running machine budgets 4.47 s for an out-of-loop takeover, so the logger keeps a 5 s ring buffer at the 1 kHz tick and freezes it at takeover. Curation searches back for the first threshold crossing. Step 2: four intervals: drift; the enforcer's blend, in which the permission path, still the only writer, ramps from the learner's command to the operator's; the operator's reorientation; and stabilized recovery from $t_{\mathrm{stab}}$ to $t_{\mathrm{rec}}$. Step 3: three rules. Mask the learner's drift commands, keep them for diagnosis. Keep the drift states for relabeling by an automated expert. Admit only the stabilized, reviewed recovery as expert labels. LINK: Lecture 5 owns intervention curation; Lecture 14 owns the causal masking rule, with the authority record and its takeover timestamps. This slide uses the mask as curation's first step. WARN: the blend and reorientation look like expert motion and are not.}
\end{frame}

\begin{frame}{Car crash recorders already keep the seconds before a trigger}
  \centering
  \small
  \renewcommand{\arraystretch}{1.25}
  \begin{tabular}{@{}lcc@{}}
    \toprule
    \textbf{Data element (49 CFR 563.7, Table I)} & \textbf{Recording interval} & \textbf{Samples/s} \\
    \midrule
    Speed, vehicle indicated & $-20$ to 0 s & 10 \\
    Accelerator pedal, \% full & $-20$ to 0 s & 10 \\
    Service brake, on/off & $-20$ to 0 s & 10 \\
    \bottomrule
  \end{tabular}\\[0.35cm]
  {\normalsize Before the phase-in: $-5.0$ to 0 s at 2 per second.\enspace The latch logger: 5 s at 1 kHz.}
  \mlsyscredit{eCFR, 49 CFR Part 563, \S563.7 Table I and footnote 4 (retrieved 2026-10-19); Physical Data, \S5.8 footnote 11.}
  \note{[1 min] [OPTIONAL] A real regulation. US rules for automotive event data recorders require the recorder to keep a fixed window of data from before the crash trigger. The current table requires speed, accelerator, and brake from 20 s before the event at 10 samples per second, 200 samples each; vehicles built before the phase-in record 5 s at 2 per second, 10 samples. The design reason matches ours. The cause of an event precedes its trigger, so a logger that starts at the trigger has already lost the drift. The latch logger keeps 5 s at the 1 kHz permission tick, 5,000 ticks, in pre-allocated memory so the freeze cannot stall. ENGAGE: why pre-allocate the buffer? (Expected: allocating memory during a safety event adds unpredictable delay.)}
\end{frame}

\begin{frame}{Expert labels on learner states make the bound linear}
  \centering
  \safeimg[width=0.90\textwidth,height=4.6cm,keepaspectratio]{bound_step3.pdf}\\[0.05cm]
  {\small \textbf{DAgger}: an expert labels the states the learner visits. Linear while $u$ stays bounded.}
  \mlsyscredit{Physical Data, \S5.8 and equation 8 (Ross et al. 2011, Theorem 3.2); $u = 1$, $S = 500$, $\varepsilon = 0.001$ (derived, conceptual).}
  \note{[2 min] [CORE] The scoreboard from the Recovery section, with one line added. Dataset Aggregation, DAgger, has an expert label the states the learner itself reaches, including curated drift states. Its bound grows as $uS\varepsilon$, where $u$ bounds how much one wrong action can raise the expert's remaining cost. With $u = 1$ at 500 steps it is 0.5, against 250 for cloning. Now the condition. If a single mistake can cost the rest of the episode, $u$ grows with $S$ and the bound returns to $S^2\varepsilon$. A bent latch pin or a sheared cable is such a state. That is the first law, irreversibility, in the language of imitation bounds. No data teaches recovery from a state with no recovery, so the permission path must keep the machine out of it. ENGAGE: is withdraw-and-reseat recoverable? Is 100 N on the latch? (Expected: yes; no.) WARN: the two bounds use different losses on different state distributions; the comparison is conceptual.}
\end{frame}

\begin{frame}{Intervention data took cable routing from 0 to 12 of 15}
  \begin{columns}[c]
    \begin{column}{0.50\textwidth}
      \centering\safeimg[width=\textwidth,height=4.6cm,keepaspectratio]{thrifty.pdf}
    \end{column}
    \begin{column}{0.46\textwidth}
      \small
      Physical robot, cable routing, 15 trials each, no interventions at test.\\[0.3cm]
      Cloning on the offline demonstrations never finished the task.\\[0.3cm]
      About 1.5 takeovers per training episode supplied the missing states.
    \end{column}
  \end{columns}
  \mlsyscredit{Hoque et al., ThriftyDAgger, CoRL 2021 (PMLR 164, 2022), Table 3 (physical cable routing; training interventions $1.55 \pm 1.16$ HG-DAgger, $1.42 \pm 1.14$ ThriftyDAgger).}
  \note{[1.5 min] [OPTIONAL] Published measurements on a real robot. Hoque and colleagues trained a cable-routing policy three ways and tested each 15 times with no human help. Behavior cloning on the offline demonstrations alone succeeded 0 of 15. Human-gated DAgger, where the supervisor decides when to take over and the takeover segments become labels, reached 10 of 15. ThriftyDAgger, where the robot asks for help when its novelty or risk estimate is high, reached 12 of 15. Both needed only about one and a half takeovers per training episode. ENGAGE: what did the takeovers add that the offline demonstrations lacked? (Expected: states the learner itself reached, with expert actions there.) WARN: 15 trials is a small sample; the gap from 0 to 10 is large, the gap from 10 to 12 is not established. Lecture 7 owns what a trial count supports.}
\end{frame}

\begin{frame}{Why does each retrained policy fail sooner?}
  \mlsysactivity{Diagnose}{1.5 min}
  \centering
  {\normalsize Each generation is retrained on all logged takeover episodes.\\Operators now intervene earlier and more often. Validation loss keeps falling.}\\[0.45cm]
  \uncover<2->{\begin{minipage}{0.86\textwidth}\small\raggedright
    \textbf{Learner-action contamination}: the drift commands entered training as expert labels, so the policy learned the moves that caused the departures.\\[0.2cm]
    \textbf{Causal confusion} differs: clean labels, but the policy keys on a consequence of the action, such as a joint current that rises because the operator already pushed.
  \end{minipage}}
  \mlsyscredit{Physical Data, \S5.8 (learner-action contamination; causal confusion, de Haan et al. 2019).}
  \note{[1.5 min] [CORE] ENGAGE: pairs, one minute. Step 2: the uncurated episodes contain the failing policy's own drift commands, labeled as if an expert issued them. The next policy learns to answer the cues of an approaching hazard with the commands that produced the hazard, drives harder toward the limit, and forces earlier takeovers. Validation loss improves because the validation split holds the same mislabeled commands. Contrast causal confusion. In de Haan and colleagues' driving example, a dashboard brake light comes on because the expert braked, so it predicts braking almost perfectly, and a policy learns to brake when the light is on. At the latch, a joint current that rises because the operator has begun to push can play that role. WARN: contamination corrupts the labels; causal confusion corrupts what the policy attends to with correct labels. Land the segment: keep the drift states, mask the drift commands, and admit only stabilized recovery.}
\end{frame}

% =============================================================================
% SECTION: PROVENANCE
% =============================================================================
\mlsysnote{[0.5 min] [CORE] Every record carries its machine with it. What must travel with the data?}
\mlsysdivider{What must travel with every recorded trajectory?}{Physical Data, \S5.9 to \S5.11}
\section{Provenance}

\begin{frame}{What does a 0.80 gripper command do on Machine B?}
  \mlsysactivity{Predict}{1.5 min}
  \centering
  \mlsyssteps[width=0.92\textwidth,height=4.8cm,keepaspectratio]{normalize_step}{2}\\[0.05cm]
  {\small\uncover<2->{Each trajectory carries its \textbf{embodiment contract}: geometry, timing, limits, compliance.}}
  \mlsyscredit{Physical Data, \S5.9.1 (normalization fallacy; derived openings).}
  \note{[1.5 min] [CORE] ENGAGE: a pipeline rescales both robots' gripper commands to the range 0 to 1 so one network head predicts both. Machine A has a 100 mm stroke; a command of 0.80 opens it to 80 mm and clears a 75 mm part. What does the same 0.80 do on Machine B, with a 60 mm stroke? Step 2: 48 mm, and the fingers land inside the part. Machine A's hold also applies 30 N, twice B's 15 N continuous rating. Every trajectory encodes its machine: sensor geometry and timing, kinematic and torque limits, compliance. The chapter calls this the embodiment contract, and borrowed actions may become labels only through a target-label gate checked against our own limit record from Lecture 2. WARN: a 7-joint arm compressed to a 6-number hand motion also loses its elbow posture; a target arm's solver may swing the elbow into a shelf. The solver works through the Jacobian, the map from joint speeds to hand speed; Lecture 11 derives it.}
\end{frame}

\begin{frame}{Open X-Embodiment pools 72 datasets with unaligned actions}
  \centering
  \small
  \renewcommand{\arraystretch}{1.2}
  \begin{tabular}{@{}llllr@{}}
    \toprule
    \textbf{Dataset} & \textbf{Robot} & \textbf{Action space} & \textbf{Rate (Hz)} & \textbf{Episodes} \\
    \midrule
    QT-Opt & Kuka iiwa & EEF position & 10 & 580,392 \\
    Language Table & xArm & EEF position & 10 & 442,226 \\
    DROID & Franka & EEF position & 15 & 92,233 \\
    ConqHose & Spot & EEF velocity & 30 & 139 \\
    MobileALOHA & MobileALOHA & Joint position & 50 & 276 \\
    \bottomrule
  \end{tabular}\\[0.3cm]
  {\small Pretraining may use them all. A target label must pass our limits.}
  \mlsyscredit{Open X-Embodiment dataset overview sheet, five of 72 rows copied (github.com/google-deepmind/open\_x\_embodiment; MLSysIM verified 2026-10-12). EEF: end effector. Physical Data, \S5.9.}
  \note{[1 min] [OPTIONAL] Five real rows of the Open X-Embodiment overview sheet. Positions, velocities, and joint angles sit side by side, at rates from 10 to 50 Hz, and the collection deliberately leaves coordinate frames unaligned. Pooling still helps; joint pretraining improved the smaller constituent datasets. The chapter separates two gates. The pretraining gate admits data for representation learning and checks provenance, frames, action semantics, and rates; the actions need not run on our robot. The target-label gate applies when an action becomes a label or a setpoint on our plant: map it through our kinematics, check our velocity and torque limits, and reject what cannot run, instead of clamping it, which would distort the motion. ENGAGE: which gate does the MobileALOHA row need before it teaches our latch? (Expected: the target-label gate, with joint-space retargeting.)}
\end{frame}

\begin{frame}{DROID's masks missed the arm, and no parser could see it}
  \begin{columns}[c]
    \begin{column}{0.58\textwidth}
      \centering\safeimg[width=\textwidth,height=4.5cm,keepaspectratio]{droid-calibration-grid.jpg}
    \end{column}
    \begin{column}{0.38\textwidth}
      \small
      The opening vote: \textbf{C}.\\[0.25cm]
      Reprojection check: project the arm from its joint angles through the stored camera pose; compare with the image.\\[0.25cm]
      A 3 mm bracket shift moves pixels; the arrays stay valid.
    \end{column}
  \end{columns}
  \mlsyscredit{Khazatsky et al., DROID, RSS 2024, Figure 13 (CC BY 4.0), via Physical Data, figure 9; \S5.10.1 and equation 9 (Lu et al. 2025).}
  \note{[1.5 min] [CORE] Back to the hook. The answer to the opening vote is C: stored camera poses that no longer matched the cameras. The check that exposes it is kinematic reprojection: take the joint angles recorded at the exposure midpoint, compute where the arm's links are, project them through the stored camera pose into the image, and measure the pixel error against the arm actually visible. A calibrated camera gives about a pixel; a bracket bumped by 3 mm or tilted by half a degree gives a large offset. The HDF5 or Zarr arrays remain valid, so a policy trained on them learns a shifted mapping from pixels to the gripper and misses targets by that offset. DROID re-estimated poses for about 36,000 episodes. ENGAGE: how often should the check run? (Expected: every session, since mounts move between sessions.) WARN: plant drift (a hot winding, a worn pad) and operator drift also change data silently; a repeated excitation trajectory each shift tests the plant.}
\end{frame}

\begin{frame}{Mars Climate Orbiter was lost to well-formed numbers}
  \begin{columns}[c]
    \begin{column}{0.40\textwidth}
      \centering\safeimg[width=\textwidth,height=4.5cm,keepaspectratio]{mars-climate-orbiter.jpg}
    \end{column}
    \begin{column}{0.56\textwidth}
      \small
      A ground file reported thruster impulse in lbf$\cdot$s; the navigation model read N$\cdot$s.\\[0.25cm]
      Every value was 4.45$\times$ too small, in range, and of the right type.\\[0.25cm]
      Planned 226 km above Mars; actual about 57 km. Lost, 1999.
    \end{column}
  \end{columns}
  \mlsyscredit{Photo: Mars Climate Orbiter during tests, NASA, public domain, Wikimedia Commons. Physical Data, War Story 5.1 (Stephenson et al. 1999).}
  \note{[1 min] [OPTIONAL] NARRATE: the Mars Climate Orbiter in test before launch. A ground software file reported thruster impulse in pound-force seconds; the trajectory models that read it assumed newton seconds, as the interface specification required. Every desaturation event entered the models too small by a factor of 4.45. The values were well formed, in range, and of the expected type, so no schema check could flag them. At orbit insertion the spacecraft passed about 57 km above the surface instead of 226 km and was lost. A robot dataset runs the same risk on every action tap and every calibration. ENGAGE: what field would have caught it? (Expected: units and frame stored with each value and checked at the interface.) FLEX: cut if short; the point lands on the next slide.}
\end{frame}

\begin{frame}{The provenance record binds an archive to its plant and gaps}
  \centering
  \footnotesize
  \renewcommand{\arraystretch}{1.18}
  \begin{tabular}{@{}>{\raggedright\arraybackslash}p{0.27\textwidth}>{\raggedright\arraybackslash}p{0.62\textwidth}@{}}
    \toprule
    \textbf{Field group} & \textbf{Latch campaign entry} \\
    \midrule
    Plant and calibration & Robot unit, firmware, pads, camera poses, clock IDs \\
    Four action taps & Asked, mapped, allowed, done, with units, frames, clamp codes \\
    Denominators & 7,353 attempts, 141.4 h, 5,000 kept, failures kept \\
    Scenario support & 1 of 6 cells; \textbf{5 known absences} \\
    Collection mode & Nominal, noise-injected, or gated intervention \\
    People recorded & Basis for recording, purpose, retention \\
    \bottomrule
  \end{tabular}\\[0.25cm]
  {\small Immutable, hashed into the container's header. New data gets a new record.}
  \mlsyscredit{Physical Data, \S5.11 and table 6; denominators from Napkin Math 5.3 (derived).}
  \note{[2.5 min] [CORE] The provenance record is the record this lecture adds to the chain, third after Lecture 2's limit record and Lecture 4's handoff record. It travels inside the data container, because format conversions lose sidecar files. It names the individual plant, because two arms of the same model differ in backlash and motor constants, so a failed attempt on a worn unit and a failed recording must be told apart per unit. It keeps all four taps with units and frames, the collection denominators, and the scenario join with its known absences. It records the people in view and the basis for recording them; Lecture 17 owns how long and for what that may be kept. latch-recovery-02 does not rewrite this record; it gets its own. LINK: Lecture 6's policy manifest cites archives by this record's digest, and Lecture 16's release manifest carries the absences. Land the segment: a stored number means nothing without its units, frame, plant, and clock, so the record binds them.}
\end{frame}

% =============================================================================
% SECTION: WRAP-UP
% =============================================================================
\mlsysnote{[0.5 min] [CORE] Check the four fallacies, then answer the hook.}
\mlsysdivider{What does an archive say about where a robot may act?}{Physical Data, \S5.12 and \S5.13}
\section{Wrap-Up}

\begin{frame}{Vote: true or false?}
  \mlsysactivity{Vote}{2 min}
  \small
  \renewcommand{\arraystretch}{1.15}
  \begin{tabular}{@{}>{\raggedright\arraybackslash}p{0.50\textwidth}>{\raggedright\arraybackslash}p{0.44\textwidth}@{}}
    \toprule
    \textbf{Claim} & \leavevmode\uncover<2->{\textbf{Verdict}} \\
    \midrule
    More samples make up for missing coverage. & \leavevmode\uncover<2->{\textbf{\textcolor{RedLine}{False.}} Repeats fill no new cells.} \\
    Dropping aborted runs cleans the dataset. & \leavevmode\uncover<2->{\textbf{\textcolor{RedLine}{False.}} It deletes the boundary states.} \\
    Labels survive a pad swap if the schema matches. & \leavevmode\uncover<2->{\textbf{\textcolor{RedLine}{False.}} The contact physics changed.} \\
    Collection scales linearly with operator hours. & \leavevmode\uncover<2->{\textbf{\textcolor{RedLine}{False.}} Wear and fatigue lower the yield.} \\
    \bottomrule
  \end{tabular}
  \mlsyscredit{Physical Data, \S5.12 (fallacies and pitfalls).}
  \note{[2 min] [CORE] ENGAGE: students write T or F for each, then compare. Step 2: all false. One: in the chapter's surface-tracking example, 72,000 more steps through the same 45 of 500 required cells leave coverage at 45. Two: in a peg-insertion pipeline, discarding the 180 of 1,000 runs that hit the 15 N cutoff leaves no states between 8 N and the trip, so a 9.5 N deployment state is unsupported. Three: stiffer fingertip pads change friction and compliance; a camera shift can be reprojected, contact cannot. Four: pilot runs on fresh hardware overstate yield; gear wear, encoder drift, and operator fatigue lower it, so budget in validated demonstrations and fatigue cycles. If many marked the second true, return to the four-zeros slide.}
\end{frame}

\begin{frame}{Back to DROID, where every file parsed}
  \begin{columns}[c]
    \begin{column}{0.42\textwidth}
      \centering\safeimg[width=\textwidth,height=3.6cm,keepaspectratio]{droid-calibration-grid.jpg}
    \end{column}
    \begin{column}{0.54\textwidth}
      \small
      Coverage: recovery states beside the corridor.\\[0.15cm]
      Labels: one tap, dated on one clock.\\[0.15cm]
      Cost: 24.6\% of an hour survives; wear is per attempt.\\[0.15cm]
      Curation: keep drift states, mask drift commands.\\[0.15cm]
      Provenance: plant, calibration, and known absences.
    \end{column}
  \end{columns}
  \mlsyscredit{Khazatsky et al., DROID, Figure 13 (CC BY 4.0). Physical Data, \S5.13.}
  \note{[1.5 min] [CORE] Return to the hook. DROID's files were valid and its camera poses were wrong; a parser checks shapes and types, and an archive's meaning lives elsewhere. Walk the five lines, one per segment. Coverage: an archive helps only where the robot's own errors take it, so it needs recovery states, audited against a table from the ODD. Labels: a label is one of four signals, dated at the exposure midpoint on one hardware clock. Cost: one scheduled hour keeps 24.6 percent as demonstrations, and every attempt wears the wrist. Curation: a takeover supplies the missing states once the learner's commands are masked. Provenance: the record binds the archive to its plant, calibration, and known absences. ENGAGE: which of these would have caught DROID's problem at collection time? (Expected: provenance with a per-session reprojection check.)}
\end{frame}

\begin{frame}{Takeaways}
  \thumbline{thumb_corridor_step2.pdf}{A policy's first error leads off the demonstrated corridor; more clean data never fills the space beside it.}\\[0.12cm]
  \thumbline{thumb_taps_step4.pdf}{A label is one of four taps, dated at the exposure midpoint on one hardware clock.}\\[0.12cm]
  \thumbline{thumb_yield_step2.pdf}{One scheduled hour keeps 24.6\% as demonstrations; wear is paid per attempt.}\\[0.12cm]
  \thumbline{thumb_scenario_step2.pdf}{Coverage is a join against a table declared from the ODD; here 1 of 6 cells.}\\[0.12cm]
  \thumbline{thumb_takeover_step3.pdf}{Takeovers supply recovery states once drift commands are masked; the provenance record keeps the absences.}
  \note{[1 min] [CORE] One line per segment. Ask students to star the one they would have gotten wrong before today. The record for the chain is the provenance record.}
\end{frame}

\begin{frame}{Open questions}
  \normalsize
  \begin{itemize}\setlength{\itemsep}{0.35cm}
    \item Robot-free hand-held grippers cost no robot wear. Which cells of the latch table can they never fill?
    \item Who should pay for an archive's known absences: more collection, a slower robot, or a narrower ODD?
    \item If a fleet logs every deployment, when does deployment data become better than demonstrations?
  \end{itemize}
  \note{[1 min] [OPTIONAL] No single answer. For the first, push toward contact forces and the target robot's own latency. A hand-held rig records no robot joints and no permission-path taps (UMI, \S5.4). For the second, the three levers are cost, speed, and scope; the permission path must refuse the absences until one of them is paid. For the third, deployment logs are censored by the safety layer and hold few boundary states (\S5.4), so they need the same curation rules. FLEX: use one as a two-minute discussion or a written assignment.}
\end{frame}

\begin{frame}{Read and explore}
  \small
  \begin{itemize}\setlength{\itemsep}{0.04cm}
    \item \textbf{Read:} Physical Data, \S5.1 to \S5.13; redo Napkin Math 5.3 with a 10 s reset.
    \item Ross, Gordon, and Bagnell, ``A Reduction of Imitation Learning and Structured Prediction to No-Regret Online Learning'' (DAgger), AISTATS 2011.
    \item Pomerleau, ``Efficient Training of Artificial Neural Networks for Autonomous Navigation,'' Neural Computation 1991.
    \item Khazatsky et al., ``DROID: A Large-Scale In-the-Wild Robot Manipulation Dataset,'' RSS 2024.
    \item Open X-Embodiment Collaboration, ``Open X-Embodiment: Robotic Learning Datasets and RT-X Models,'' ICRA 2024.
    \item Gebru et al., ``Datasheets for Datasets,'' Communications of the ACM 2021.
  \end{itemize}
  \note{[0.5 min] [CORE] The chapter first. Each reading backs one idea: DAgger and the recoverability constant (Ross et al.), the recovery gap and its first fix (Pomerleau), a single-embodiment corpus and its recalibration (DROID), pooling across embodiments (Open X-Embodiment), and documentation for static corpora that the provenance record extends (Gebru et al.). In the redo, a 10 s reset raises attempts per session; ask which bill moves.}
\end{frame}

\begin{frame}[c]
  \centering
  {\small\textcolor{lightgray}{NEXT LECTURE \textbar{} POLICY TRAINING}}\\[0.4cm]
  \includegraphics[width=3.2cm,height=1.4cm,keepaspectratio]{thumb_bound_step2.pdf}\\[0.35cm]
  {\LARGE\bfseries Why can a policy improve its training score\\while becoming less reliable on the machine?}\\[0.3cm]
  {\normalsize\textcolor{midgray}{The latch archive supports 1 of 6 cells; training will not show where the other 5 begin.}}
  \note{[0.5 min] [CORE] Close on the archive's known absences. Lecture 6 trains two latch candidates on this archive, one cloned from the demonstrations and one distilled from a simulator that learned the fast approach no one collected, and writes what each may claim into a policy manifest whose domain cannot exceed the cells this provenance record supports.}
\end{frame}

\end{document}
